\section{Gaussian truncation with a bounded local neighbourhood}
\label{sec:hastings05_metric_filter}

Lemma~\ref{thm:hastings05_metric_filter} constructs the local Gaussian
filters of~\cite[eq.~(mid), lines 285--349]{Hastings2006Solving} on a finite
metric space. A bound on the number of sites in a ball of the interaction
range suffices. Polynomial growth of larger spheres is unnecessary
\cite{gap:hastings05_geometric_growth_scope}.

\begin{lemma}[An exponentially weighted graph tail]
    \label{lem:hastings05_degree_tail}
    \lean{TNLean.Lattice.graphClosedBall,
        TNLean.Lattice.card_graphClosedBall_le,
        TNLean.Lattice.sum_exp_neg_edist_le_degree,
        TNLean.Lattice.sum_weight_exp_neg_edist_tail_le_degree}
    \leanok
    Let $G$ be a finite graph of maximum degree at most $D$. Write
    $d_G(x,y)=\infty$ for vertices in different components, and interpret
    $e^{-\mu\infty}$ as zero. If $\mu>\log(D+1)$, then
    \begin{align}
        \sum_y e^{-\mu d_G(x,y)}
        &\le \frac{1}{1-e^{-(\mu-\log(D+1))}}.
        \label{eq:hastings05_degree_sum}
    \end{align}
    If each vertex $y$ carries nonnegative weights with total at most $B$,
    the weighted sum over anchors at distance at least $a$ is at most
    $B e^{-\gamma a}/(1-e^{-\gamma})$, where
    $\gamma=(\mu-\log(D+1))/2$.
\end{lemma}
\begin{proof}\leanok
    A graph ball of integer radius $k$ has at most $(D+1)^k$ vertices.
    Sum over distance spheres. For the tail, split the excess decay
    $\mu-\log(D+1)$ equally between the distance threshold and the
    remaining geometric series.
\end{proof}

\begin{lemma}[A termwise truncation estimate]
    \label{lem:hastings05_weighted_truncation}
    \lean{TNLean.Lattice.norm_gaussianFilteredTerm_sub_truncated_le_of_commutator_bounds,
        TNLean.Lattice.norm_commutator_heisenbergEvolution_le_anchor_kernel,
        TNLean.Lattice.norm_gaussianFilteredTerm_sub_truncated_le_anchor_sum,
        TNLean.Lattice.norm_gaussianFilteredTerm_sub_truncated_le_degree_tail}
    \leanok
    \uses{lem:hastings05_degree_tail}
    Let $h_k=h_k^\dagger$, let $H=\sum_k h_k$ and
    $H_A=\sum_{k\in A}h_k$, and let $W$ be an
    operator. Suppose $T\ge0$ and
    $\|[h_k,e^{itH_A}We^{-itH_A}]\|\le C_k$ for $k\notin A$ and
    $|t|\le T$. For the normalized Gaussian filter $\mathcal F_{q,\Delta}$,
    with $q\ge0$ and the usual totalized convention at $q=0$,
    \begin{align}
        \|\mathcal F_{q,\Delta}^{H}(W)
            -\mathcal F_{q,\Delta}^{H_A}(W)\|
        &\le T\sum_{k\notin A}C_k
            +2\|W\|e^{-(T\Delta)^2/(2q)}.
        \label{eq:hastings05_weighted_truncation}
    \end{align}
    For local interactions on a bounded-degree graph, the omitted sum is
    bounded by Lemma~\ref{lem:hastings05_degree_tail} and the
    Lieb--Robinson estimate at a freely chosen rate
    $\mu>\log(D+1)$.
\end{lemma}
\begin{proof}\leanok
    Apply the Duhamel comparison to the two evolutions, retaining the
    individual omitted commutators before taking their sum. Split the
    Gaussian integral at $|t|=T$. The integral outside this interval is
    bounded by its Gaussian tail. For the graph estimate, use the finite
    support propagation bound at each omitted anchor and sum its
    exponentially decaying weights.
\end{proof}

\begin{lemma}[The actual metric-ball Gaussian filters]
    \label{thm:hastings05_metric_filter}
    \lean{TNLean.Lattice.metricRangeGraph,TNLean.Lattice.metricRangeBall,
        TNLean.Lattice.exists_metric_truncated_filter_budget,
        TNLean.Lattice.metric_truncated_filter_properties,
        TNLean.Lattice.exists_metric_truncated_filter_localization,
        TNLean.Lattice.exists_metric_truncated_filter_sum_bound}
    \leanok
    \uses{lem:hastings05_weighted_truncation}
    Fix $R>0$, $J\ge0$, $\Delta>0$, and an integer $S$. There are
    $c_1,\eta>0$ and $C\ge0$, depending only on these parameters, with
    the following property. Let $\Lambda$ be a finite metric space, let
    $H=\sum_i h_i$, and suppose $h_i=h_i^\dagger$ is supported in
    $B_R(i)$, $\|h_i\|\le J$, and $|B_R(i)|\le S$. For $\ell\ge R$, set
    \begin{align}
        q&=c_1\ell\Delta,&
        H_i^{\mathrm{loc}}&=\sum_{j:\,d(i,j)\le\ell-R}h_j,&
        M_i&=\mathcal F_{q,\Delta}^{H_i^{\mathrm{loc}}}(h_i).
        \label{eq:hastings05_metric_filter}
    \end{align}
    Then $M_i=M_i^\dagger$ is supported in $B_\ell(i)$,
    $\|M_i\|\le J$, and
    \begin{align}
        \|h_i\|e^{-q/2}
          +\|\mathcal F_{q,\Delta}^{H}(h_i)-M_i\|
        &\le \|h_i\|C e^{-\eta\ell},
        \label{eq:hastings05_metric_budget}\\
        \Big\|\sum_i M_i-H\Big\|
        &\le |\Lambda|JC e^{-\eta\ell}.
        \label{eq:hastings05_metric_sum}
    \end{align}
    If the spectrum of $H$ is contained in
    $\{E_0\}\cup[E_0+\Delta,\infty)$, the component of $M_i\psi$
    outside the full ground space is at most
    $JC e^{-\eta\ell}\|\psi\|$ for every ground vector $\psi$.
    If the ground space is a line spanned by a unit vector $\Omega$, this
    gives $\|(M_i-\langle\Omega,M_i\Omega\rangle)\Omega\|
    \le JC e^{-\eta\ell}$.
\end{lemma}
\begin{proof}\leanok
    Join two sites when their metric distance is at most $R$. This graph
    has degree at most $S$, and each interaction has graph diameter at
    most two. Choose $\mu=\log(S+1)+1$ in
    Lemma~\ref{lem:hastings05_weighted_truncation}. Choose $T=c_1\ell$
    with $c_1>0$ sufficiently small that the spatial decay dominates the
    time growth. The Gaussian tail is $e^{-q/2}$ because $q=T\Delta$.
    A fixed fraction of the remaining exponential decay absorbs the
    factor $\ell$. Local support follows from the triangle inequality.
    The full filters sum exactly to $H$, which proves
    (\ref{eq:hastings05_metric_sum}). The ground-space conclusions
    follow from the spectral suppression of the full Gaussian filter.
\end{proof}

\begin{lemma}[Zero interaction range]
    \label{lem:hastings05_zero_range}
    \lean{TNLean.Lattice.gaussianFilteredTerm_eq_of_commute,
        TNLean.Lattice.commute_of_singleton_support,
        TNLean.Lattice.gaussianFilteredTerm_full_and_sum_eq_of_zero_range,
        TNLean.Lattice.starProjection_term_eq_of_zero_range,
        TNLean.Lattice.sub_starProjection_gaussianFilteredTerm_sum_eq_zero_of_zero_range}
    \leanok
    When $R=0$ on a metric space, the interactions have single-site
    support and commute. Every full or truncated Gaussian filter of
    $h_i$ equals $h_i$. Each term preserves the full ground eigenspace.
\end{lemma}
\begin{proof}\leanok
    Operators supported at different single sites commute. Conjugating
    $h_i$ by any sum of these operators leaves it unchanged. Commutation
    with $H$ also preserves its eigenspaces.
\end{proof}
